% !TEX root = ../main.tex
\clearpage
\section{Criterios de adecuación al hospital}
\label{sec:evaluacion}
La revisión documental delimita candidatos y preguntas pendientes. Los criterios siguientes orientan la selección; \textbf{el diseño del piloto, las métricas y los umbrales de aceptación siguen pendientes de definición}.

\subsection{Protección de datos, seguridad e integración}
Se propone priorizar \textbf{almacenamiento y procesamiento en la UE/EEE, ENS de categoría ALTA e integración con Dedalus HCIS}. La ubicación europea no acredita por sí sola cumplimiento del RGPD: deben verificarse finalidades, responsabilidades, contratos, subencargados y transferencias~\cite{rgpd2016}. Para ENS deben comprobarse titular, sistema, alcance, emisor y vigencia, así como la categoría exigible al sistema concreto~\cite{ens2022}.

La arquitectura dependerá del tratamiento de datos admisible y la capacidad operativa del hospital. Las certificaciones no sustituyen la evaluación clínica; una alianza con Dedalus o una API tampoco acredita integración con la instalación local de HCIS (apéndices~\ref{app:certificaciones} y~\ref{app:hcis}).

\subsection{Factores de comparación}
\begin{small}
\begin{longtable}{L{.22\linewidth}L{.70\linewidth}}
\toprule
Factor & Evidencia relevante para la decisión \\
\midrule\endhead
Seguridad y datos & Certificados e informes completos; alojamiento y procesamiento; acceso de soporte; retención de audio y notas; subencargados y reutilización para entrenamiento. \\
Calidad en español & Evaluaciones en el idioma y contexto previstos: omisiones, adiciones, negaciones, fármacos, dosis y atribución de hablantes. Resultados por especialidad y condiciones de captura. \\
Utilidad clínica & Tiempo total de documentación, carga de revisión, experiencia del profesional y del paciente, utilización y abandono. \\
Integración y operación & Identificación de paciente y episodio, edición, transferencia, firma, registro de errores, dispositivos compatibles y comportamiento ante pérdida de conexión. \\
Coste total & Licencias, consumo, equipos, integración, mantenimiento, formación, soporte, mínimos contractuales y condiciones de salida. \\
\bottomrule
\end{longtable}
\end{small}
La matriz combina las necesidades del proyecto con dominios de estudios comparativos~\cite{NNWFT2JX,PRTISMYF}. Declarar capacidad multilingüe no acredita calidad equivalente: interesa evaluar español, acentos, expresiones coloquiales y varios hablantes. Las simulaciones disponibles orientan las pruebas, sin sustituir la validación local.

\subsection{Interpretación de la evidencia}
Se distinguen documentos examinados, declaraciones comerciales, evidencia de uso y aspectos pendientes. Revisión por pares, independencia del proveedor y aplicabilidad local se valoran por separado. Los estudios de implantación, los ensayos comparativos y las evaluaciones de errores responden a preguntas diferentes; sus diseños y medidas desaconsejan una puntuación global de superioridad clínica~\cite{956B6PZT}.

La similitud textual no equivale a exactitud clínica ni la nota del médico constituye una referencia infalible. La evaluación distinguirá calidad del borrador, correcciones y calidad de la nota final~\cite{4IHZ7QFQ,ZI87XIWC}.
